\section{Additional Experiments on Comparing Various Text Moderation Models} 

To quantitatively assess the efficacy of the current moderation systems, we conducted a few experiments on two publicly available text moderation systems that are widely adopted: OpenAI Moderation API and Persepctive API. We prompted these moderation systems with the same evaluation dataset that we used in producing Figure~\ref{fig:flagged_proportion}, and we used this data to measure the agreement between the underlying moderation system and those three external evaluators presented in Figure~\ref{fig:flagged_proportion}. We fed the system with Q and A concatenated. The experiment result is best presented in the worksheet format, so it is provided \textbf{in the supplementary materials}. From the results, we have concluded a few things:

\textbf{Perspective API:}
\begin{itemize}[leftmargin=*]
    \item Its ability to comprehend context appears limited, as evidenced by the consistently low harm scores assigned to responses from Alpaca-7B and Alpaca-13B in the categories of "Terrorism, Organized Crime," "Animal Abuse," "Non-Violent Unethical Behavior," and "Drug Abuse, Weapons, Banned Substance." In these cases, humans, our QA moderation, and GPT-4 (collectively referred to as the three evaluators) all agreed that the response was harmful.
    \item It is highly sensitive to specific keywords. It's important to note that all prompts in our evaluation dataset are malicious, and some may contain explicit language. Despite this, GPT-3.5 emerges as the safest model, with nearly all of its responses being rated non-harmful by the three evaluators. However, Perspective API still flags texts as harmful, regardless of the response's appropriateness. This trend is apparent in the "gpt-3.5-turbo" and "vicuna-7b" responses within the "Sexually Explicit, Adult Content" category.
    \item The effectiveness of the API's detection, measured in terms of harm category probability, correlates strongly with text length. The presence of additional text without harmful keywords tends to dilute the output probability.
\end{itemize}


\textbf{OpenAI Moderation API:}
\begin{itemize}[leftmargin=*]
    \item OpenAI showed signs of context comprehension, as indicated by the decreasing trend in the proportion of flagged responses from Alpaca-7B 
 Alpaca-13B >> Vicuna-7B > gpt-3.5-turbo - with a lower proportion being better. This trend is consistent with the findings observed in the evaluation results provided by the three evaluators.
    \item However, the issue with OpenAI Moderation is its lack of robustness against unsafe QA pairs. Due to its smaller scope in the categorization of harms, it has failed to identify any harms in several categories, 'Terrorism, Organized Crime', 'Privacy Violation', 'Drug Abuse, Weapons, Banned Substance', in Alpaca-7B and Alpaca-13B responses. Notably in these categories, at least 50\% or more of the QA pairs were unanimously flagged as harmful by all three evaluators.
    \item OpenAI Moderation can also exhibit over-sensitivity towards safe QA pairs. For instance, in the category "Sexually Explicit, Adult Content" for responses generated by gpt-3.5-turbo and Vicuna-7B, OpenAI Moderation marked 10\% and 20\% of the QA pairs as unsafe, respectively. However, all pairs in these cases were unanimously deemed safe by all three evaluators.
\end{itemize}

Based on these findings, we have determined that the Perspective API is not suited for the QA moderation task. Its reliance on keyword detection means that any chatbot system using this moderation approach is likely to experience frequent request rejections, as the API's hypersensitivity to certain keywords may trigger false alarms as opposed to addressing the user's request. While OpenAI moderation demonstrates some capability of performing the QA moderation task, it is not robust enough when compared to our moderation model.